% ID: SEL-D0081-OB-E01
% Tanda S081.
\textbf{a)} Despejamos las restricciones:
\[
 y\le7x+10,\qquad y\le2-x,\qquad y\ge\frac{3x-14}{5}.
\]
El recinto es el triángulo cerrado que queda por debajo de las dos primeras rectas y por encima de la tercera. Intersecando las fronteras por parejas:
\[
 \begin{array}{c|c}
 \text{Rectas}&\text{Vértice}\\\hline
 y=7x+10,\ y=2-x&(-1,3)\\
 y=7x+10,\ 5y=3x-14&(-2,-4)\\
 y=2-x,\ 5y=3x-14&(3,-1)
 \end{array}
\]
Cada vértice satisface la restricción restante.
\begin{center}
\setlength{\unitlength}{1mm}
\begin{picture}(120,100)(-45,-50)
\put(-40,0){\vector(1,0){106}}\put(0,-45){\vector(0,1){86}}
\put(68,0){$x$}\put(1,41){$y$}
{\color{gray}
\qbezier(-33,-43.2)(-22.5,-4)(-12,35.2)
\qbezier(-37.5,36)(11.25,10)(60,-16)
\qbezier(-37.5,-34.4)(11.25,-18.8)(60,-3.2)}
\linethickness{0.45mm}
\qbezier(-15,24)(-22.5,-4)(-30,-32)
\qbezier(-30,-32)(7.5,-20)(45,-8)
\qbezier(45,-8)(15,8)(-15,24)
\put(-26,28){\scriptsize $(-1,3)$}
\put(-40,-40){\scriptsize $(-2,-4)$}
\put(46,-12){\scriptsize $(3,-1)$}
\end{picture}
\end{center}

\textbf{b)} Como la función es lineal y el recinto es un polígono cerrado y acotado, comparamos los vértices:
\[
 F(-1,3)=7,\qquad F(-2,-4)=-16,\qquad F(3,-1)=3.
\]
El máximo es $\boxed{7}$ en $(-1,3)$ y el mínimo es $\boxed{-16}$ en $(-2,-4)$. Ambos se alcanzan en un único punto.
